% Original CSC302 teaching exposition based on Macbeth, Chapter 3, Sections 3.1--3.5.
% Source: https://hrmacbeth.github.io/math2001/03_Parity_and_Divisibility.html
% Source snapshot: hrmacbeth/math2001 e660f42b13ddcb6d12b52ba036d6bd071a0cfb9b.
% All prose is independently written. New numerical variants use the cited proof method.
% CSC302 lecture slides adapted from Heather Macbeth, The Mechanics of Proof.
\documentclass[aspectratio=169,11pt]{beamer}
\usepackage[T1]{fontenc}
\usepackage{lmodern}
\usepackage{amsmath,amssymb,mathtools}
\usepackage{booktabs,array,listings,tikz}
\usetikzlibrary{arrows.meta,positioning,calc}
\definecolor{ProofNavy}{HTML}{18324F}
\definecolor{ProofTeal}{HTML}{007D82}
\definecolor{ProofPaper}{HTML}{F3F6F8}
\usetheme{default}
\setbeamercolor{normal text}{fg=ProofNavy,bg=white}
\setbeamercolor{structure}{fg=ProofTeal}
\setbeamercolor{frametitle}{fg=ProofNavy}
\setbeamercolor{block title}{fg=white,bg=ProofNavy}
\setbeamercolor{block body}{fg=ProofNavy,bg=ProofPaper}
\setbeamercolor{block title example}{fg=white,bg=ProofTeal}
\setbeamercolor{block body example}{fg=ProofNavy,bg=ProofPaper}
\setbeamerfont{frametitle}{size=\Large,series=\bfseries}
\setbeamertemplate{navigation symbols}{}
\setbeamertemplate{itemize item}{\small\textbullet}
\setbeamertemplate{footline}{%
  \leavevmode\hbox{%
  \begin{beamercolorbox}[wd=.86\paperwidth,ht=2.8ex,dp=1.2ex,leftskip=1.6em]{author in head/foot}%
    \scriptsize CSC302\quad\textbar\quad\insertshorttitle
  \end{beamercolorbox}%
  \begin{beamercolorbox}[wd=.14\paperwidth,ht=2.8ex,dp=1.2ex,rightskip=1.6em plus1fil]{date in head/foot}%
    \scriptsize\insertframenumber/\inserttotalframenumber
  \end{beamercolorbox}}}
\setbeamersize{text margin left=8mm,text margin right=8mm}
\lstdefinelanguage{Lean}{%
  morekeywords={example,theorem,lemma,def,by,calc,have,exact,apply,intro,intros,constructor,rcases,rfl,rw,ring,numbers,linarith,nlinarith,rel,extra,addarith,obtain,left,right,use,choose,dsimp,change,ext,fun,show,from,at,with,induction,case,decide,contradiction,push_neg,by_cases,by_contra,assumption,existsi,simp,simpa},
  sensitive=true,morecomment=[l]{--},morestring=[b]"}
\lstset{language=Lean,basicstyle=\ttfamily\small,keywordstyle=\color{ProofTeal}\bfseries,
  commentstyle=\color{gray},stringstyle=\color{ProofNavy},showstringspaces=false,
  columns=fullflexible,keepspaces=true,breaklines=true,frame=single,
  rulecolor=\color{ProofTeal!30},backgroundcolor=\color{ProofPaper},
  aboveskip=0.5em,belowskip=0.4em}
\newcommand{\N}{\mathbb{N}}
\newcommand{\Z}{\mathbb{Z}}
\newcommand{\Q}{\mathbb{Q}}
\newcommand{\R}{\mathbb{R}}
\newcommand{\coursechapter}[2]{%
  \title[Chapter #1: #2]{Chapter #1\\#2}%
  \subtitle{The Mechanics of Proof \textbar{} CSC302 lecture notes}%
  \author{Based on Heather Macbeth's textbook}%
  \date{}%
}
\newcommand{\sourcecredit}[1]{\par\vfill{\tiny Source: Heather Macbeth, \emph{The Mechanics of Proof}, #1.}}
\newcommand{\bookurl}{\url{https://hrmacbeth.github.io/math2001/}}
\hypersetup{colorlinks=true,urlcolor=ProofTeal,linkcolor=ProofNavy}

\coursechapter{3}{Parity and divisibility}
\begin{document}

\begin{frame}\titlepage\end{frame}

\begin{frame}{What you will learn}
\begin{itemize}
\item Unfold evenness, oddness, divisibility, and congruence into witness statements.
\item Construct new witnesses by algebraic calculation.
\item Prove and use rules for modular arithmetic.
\item Use residue cases to cover all integers.
\item Explain how a B\'ezout identity justifies certain cancellation steps.
\end{itemize}
\begin{block}{The recurring proof pattern}
Unpack the hypothesis $\longrightarrow$ calculate $\longrightarrow$ identify a witness for the conclusion.
\end{block}
\end{frame}

\section{3.1 Definitions; parity}
% Definitions and witness methods: Examples 3.1.1--3.1.7.
\begin{frame}{3.1\quad Definitions tell us what must be proved}
For an integer $n$:
\begin{align*}
n\text{ is even}&\quad\Longleftrightarrow\quad \exists k\in\Z,\ n=2k,\\
n\text{ is odd}&\quad\Longleftrightarrow\quad \exists k\in\Z,\ n=2k+1.
\end{align*}
\begin{itemize}
\item To prove evenness or oddness, produce an integer $k$.
\item To use evenness or oddness, introduce the promised $k$ and its equality.
\item Different existential statements may have different witnesses, even if both are written with the letter $k$.
\end{itemize}
Lean: \texttt{dsimp [Int.Odd]} displays the textbook's definition.
\sourcecredit{Section 3.1}
\end{frame}

\begin{frame}{Worked examples: zero and negative integers count too}
\begin{center}
\begin{tabular}{lll}
\toprule
Claim & Witness & Verification \\
\midrule
$0$ is even & $k=0$ & $0=2\cdot0$ \\
$-8$ is even & $k=-4$ & $-8=2(-4)$ \\
$9$ is odd & $k=4$ & $9=2\cdot4+1$ \\
$-7$ is odd & $k=-4$ & $-7=2(-4)+1$ \\
\bottomrule
\end{tabular}
\end{center}
\begin{block}{Why the domain matters}
The witnesses for $-8$ and $-7$ are integers, not natural numbers.
Parity is not restricted to positive integers.
\end{block}
\sourcecredit{Section 3.1; methods of Examples 3.1.1--3.1.2}
\end{frame}

\begin{frame}[fragile]{Worked example: turn one oddness witness into another}
If $n=2k+1$, then
\[
5n-2=5(2k+1)-2=2(5k+1)+1.
\]
Thus $5n-2$ is odd, with witness $5k+1$.
\begin{lstlisting}
example {n : Int} (h : Int.Odd n) : Int.Odd (5 * n - 2) := by
  dsimp [Int.Odd] at *
  cases h with
  | intro k hk =>
      use 5 * k + 1
      calc
        5 * n - 2 = 5 * (2 * k + 1) - 2 := by rw [hk]
        _ = 2 * (5 * k + 1) + 1 := by ring
\end{lstlisting}
\sourcecredit{Section 3.1; method of Example 3.1.3}
\end{frame}

% Parity case split: method of Example 3.1.9, with the constant 4 removed.
\begin{frame}{Worked example: every integer has one of two parities}
Every integer is even or odd. Use this fact to prove $n^2+n$ is even.
\begin{block}{Case 1: $n=2k$}
\[
n^2+n=(2k)^2+2k=2(2k^2+k).
\]
\end{block}
\begin{block}{Case 2: $n=2k+1$}
\[
n^2+n=(2k+1)^2+(2k+1)=2(2k^2+3k+1).
\]
\end{block}
Both witnesses are integers, and the cases exhaust all integers.
Lean supplies the case split as \texttt{Int.even\_or\_odd n}.
\sourcecredit{Example 3.1.9; parity dichotomy revisited in Chapter 4}
\end{frame}

\begin{frame}{Practice 3.1: name the output witness}
\begin{enumerate}
\item If $n$ is even, prove that $n^2+3$ is odd.
\item If $m,n$ are odd integers, prove that $mn$ is odd.
\end{enumerate}
\vspace{1em}
Use different letters for witnesses from independent hypotheses.
End each proof by explicitly identifying an integer $k$ in the form $2k+1$.
\end{frame}

\begin{frame}{Practice 3.1: solution discussion}
\textbf{1.} Write $n=2r$ for some $r\in\Z$. Then
\[
n^2+3=4r^2+3=2(2r^2+1)+1.
\]
The output witness is $2r^2+1$.
\medskip

\textbf{2.} Write $m=2a+1$ and $n=2b+1$, with $a,b\in\Z$. Then
\[
mn=(2a+1)(2b+1)=2(2ab+a+b)+1.
\]
The output witness is $2ab+a+b$.
\sourcecredit{Problem 2 is a standard parity fact also practiced in Section 3.1}
\end{frame}

\section{3.2 Divisibility}
\begin{frame}{3.2\quad Divisibility is an existence statement}
For $a,b\in\Z$, define
\[
a\mid b\quad\Longleftrightarrow\quad\exists k\in\Z,\ b=ak.
\]
We say that $a$ divides $b$, that $a$ is a factor of $b$, or that $b$ is a multiple of $a$.
\begin{itemize}
\item $6\mid42$, because $42=6\cdot7$.
\item $-3\mid12$, because $12=(-3)(-4)$.
\item $a\mid0$ for every $a$, because $0=a\cdot0$.
\item $0\mid b$ forces $b=0$.
\end{itemize}
For divisibility in $\N$, the witness must lie in $\N$.
\sourcecredit{Examples 3.2.1--3.2.2 and Section 3.2 exercises}
\end{frame}

% Witness transport: Example 3.2.3.
\begin{frame}[fragile]{Worked example: retain a common factor}
If $a\mid b$, write $b=ak$. Then
\[
b^2+2b=(ak)^2+2ak=a\bigl(k(ak+2)\bigr).
\]
This proves $a\mid b^2+2b$.
\begin{lstlisting}
example {a b : Int} (h : Dvd.dvd a b) :
    Dvd.dvd a (b ^ 2 + 2 * b) := by
  cases h with
  | intro k hk =>
      use k * (a * k + 2)
      calc
        b ^ 2 + 2 * b = (a * k) ^ 2 + 2 * (a * k) := by rw [hk]
        _ = a * (k * (a * k + 2)) := by ring
\end{lstlisting}
\sourcecredit{Example 3.2.3; ASCII form of the divisibility notation}
\end{frame}

\begin{frame}{Two basic facts and the hypotheses they need}
\textbf{Transitivity.} For integers, $a\mid b$ and $b\mid c$ imply $a\mid c$.
Indeed, if $b=ar$ and $c=bs$, then
\[
c=(ar)s=a(rs).
\]
\begin{block}{A size bound in $\N$}
If $a,b\in\N$, $b>0$, and $a\mid b$, then $0<a\le b$.
\end{block}
Write $b=ak$. Since $ak>0$, both $a$ and $k$ are positive.
Thus $k\ge1$, and $a=a\cdot1\le ak=b$.
\medskip

The condition $b>0$ is essential: $8\mid0$, but $8\nleq0$.
\sourcecredit{Section 3.2; witness composition and Examples 3.2.7--3.2.8}
\end{frame}

\begin{frame}{Worked example: prove non-divisibility by a gap}
\begin{block}{Consecutive multiples test}
Let $d>0$ and $q,t\in\Z$. If $dq<t<d(q+1)$, then $d\nmid t$.
\end{block}
Why? If $t=dk$ for an integer $k$, dividing by positive $d$ gives
$q<k<q+1$, impossible for an integer.
\medskip

\textbf{Example.} To show $4\nmid-6$, observe
\[
4(-2)=-8<-6<-4=4(-2+1).
\]
The same test works for negative $t$; choose the two surrounding multiples.
\sourcecredit{Section 3.2; method of Example 3.2.6}
\end{frame}

\begin{frame}{Practice 3.2: compose the witnesses}
\begin{enumerate}
\item For $a,b\in\Z$, prove that $a\mid b$ implies $a\mid3b-b^2$.
\item For $m,n,p\in\Z$, prove that $m\mid n$ and $n^3\mid p$ imply $m^3\mid p$.
\item Show that $5\nmid17$ by placing $17$ between consecutive multiples.
\end{enumerate}
Do not divide by $a$ or $m$: those numbers may be zero.
Use the equalities supplied by divisibility instead.
\end{frame}

\begin{frame}{Practice 3.2: solution discussion}
\textbf{1.} If $b=ak$, then
\[
3b-b^2=3ak-a^2k^2=a(3k-ak^2).
\]
\textbf{2.} Write $n=mr$ and $p=n^3s$. Then
\[
p=(mr)^3s=m^3(r^3s).
\]
\textbf{3.} Since $5\cdot3=15<17<20=5(3+1)$, the number $17$ is not a multiple of $5$.
\end{frame}

\section{3.3 Modular arithmetic: theory}
\begin{frame}{3.3\quad Congruence compares a difference to the modulus}
For integers $a,b,m$, define
\[
a\equiv b\pmod m\quad\Longleftrightarrow\quad m\mid(a-b).
\]
Equivalently, $a-b=mk$ for some $k\in\Z$.
\begin{exampleblock}{Worked example}
$-8\equiv2\pmod5$, because $-8-2=5(-2)$.
\end{exampleblock}
\begin{itemize}
\item Congruent numbers need not be equal.
\item For $m>0$, congruence means having the same remainder modulo $m$.
\item The definition also makes sense for $m=0$: it reduces to $a=b$.
\end{itemize}
Lean's textbook notation is represented by \texttt{Int.ModEq m a b}.
\sourcecredit{Section 3.3 definition and Examples 3.3.1--3.3.2}
\end{frame}

% Theorem and proof methods: Examples 3.3.3--3.3.5.
\begin{frame}{Addition and subtraction preserve congruence}
Assume $a\equiv b\pmod m$ and $c\equiv d\pmod m$.
Choose integers $r,s$ with $a-b=mr$ and $c-d=ms$.
Then
\[
\begin{aligned}
(a+c)-(b+d)&=(a-b)+(c-d)\\
           &=mr+ms=m(r+s),\\[.4em]
(a-c)-(b-d)&=(a-b)-(c-d)\\
           &=mr-ms=m(r-s).
\end{aligned}
\]
Therefore $a+c\equiv b+d\pmod m$ and $a-c\equiv b-d\pmod m$.
Negation is similar: $(-a)-(-b)=-(a-b)=m(-r)$.
\sourcecredit{Examples 3.3.3--3.3.5}
\end{frame}

% Theorem and proof method: Example 3.3.6.
\begin{frame}{Multiplication preserves congruence}
Again assume $a-b=mr$ and $c-d=ms$. The useful identity is
\[
ac-bd=(a-b)c+b(c-d).
\]
It gives
\[
\begin{aligned}
ac-bd&=(mr)c+b(ms)\\
     &=m(rc+bs).
\end{aligned}
\]
Hence $ac\equiv bd\pmod m$, with witness $rc+bs$.
\begin{block}{A useful special case}
$a^2-b^2=(a-b)(a+b)=m\bigl(r(a+b)\bigr)$, so $a^2\equiv b^2\pmod m$.
\end{block}
\sourcecredit{Examples 3.3.6 and 3.3.8}
\end{frame}

\begin{frame}{Congruence supports chains and powers}
For integers $a,b,c,m$:
\begin{itemize}
\item \textbf{Reflexivity:} $a-a=m\cdot0$, so $a\equiv a\pmod m$.
\item \textbf{Symmetry:} if $a-b=mr$, then $b-a=m(-r)$.
\item \textbf{Transitivity:} if $a-b=mr$ and $b-c=ms$, then $a-c=m(r+s)$.
\end{itemize}
\begin{block}{Power rule}
If $a\equiv b\pmod m$ and $k\in\N$, then $a^k\equiv b^k\pmod m$.
\end{block}
For $k=0$ both sides are $1$. Repeated multiplication proves each positive-power case;
Chapter 6 organizes this argument as induction on $k$.
\sourcecredit{Examples 3.3.8--3.3.10, Section 3.3 exercises; Example 6.1.3}
\end{frame}

\begin{frame}{Ordinary division is not a congruence rule}
Modulo $6$ we have
\[
2\cdot1\equiv2\cdot4\pmod6,
\]
because $2-8=-6$ is divisible by $6$.
But cancelling the $2$ would assert
\[
1\equiv4\pmod6,
\]
which is false: $1-4=-3$ is not divisible by $6$.
\begin{block}{What failed?}
An equality can be divided by a nonzero real factor.
A congruence needs additional arithmetic information about the factor and modulus.
\end{block}
Section 3.5 supplies identities that make certain cancellations valid.
\sourcecredit{Section 3.3; new counterexample illustrating Example 3.3.7}
\end{frame}

\begin{frame}{Practice 3.3: prove the rule from the definition}
Suppose $a\equiv b\pmod m$.
\begin{enumerate}
\item Prove $2a+1\equiv2b+1\pmod m$ by giving a witness.
\item Prove $a^3\equiv b^3\pmod m$ using a factorization.
\end{enumerate}
\vspace{1em}
Start from $a-b=mr$.
Your calculation should end with $m$ times an integer.
\end{frame}

\begin{frame}{Practice 3.3: solution discussion}
Choose $r\in\Z$ with $a-b=mr$.
\textbf{1.}
\[
(2a+1)-(2b+1)=2(a-b)=2mr=m(2r).
\]
\textbf{2.}
\[
\begin{aligned}
a^3-b^3&=(a-b)(a^2+ab+b^2)\\
       &=m\bigl(r(a^2+ab+b^2)\bigr).
\end{aligned}
\]
The witnesses are $2r$ and $r(a^2+ab+b^2)$, respectively.
\sourcecredit{Problem 2 follows Example 3.3.9}
\end{frame}

\section{3.4 Modular arithmetic: calculations}
\begin{frame}{3.4\quad Calculate with the rules you have proved}
Let $a,b\in\Z$, with $a\equiv3\pmod7$ and $b\equiv4\pmod7$.
Find the residue of $a^2+b$.
\[
\begin{aligned}
a^2+b&\equiv3^2+4\pmod7&&\text{power and addition rules}\\
     &=13&&\text{arithmetic}\\
     &\equiv6\pmod7&&13-6=7\cdot1.
\end{aligned}
\]
\begin{block}{Why this is a proof}
Each congruence uses an established rule; an equality also gives a congruence.
Transitivity connects the initial expression to its final residue.
\end{block}
\sourcecredit{Section 3.4; method of Example 3.4.2}
\end{frame}

% Source theorem and case method: Example 3.4.4.
\begin{frame}{Worked example: cover every possible residue}
\textbf{Claim.} Every integer $x$ satisfies $x^3\equiv x\pmod3$.

There are three possible residues of $x$ modulo $3$:
\begin{center}
\begin{tabular}{cccc}
\toprule
Residue $r$ & $r^3$ & $r^3-r$ & Divisible by $3$? \\
\midrule
$0$ & $0$ & $0$ & $0=3\cdot0$ \\
$1$ & $1$ & $0$ & $0=3\cdot0$ \\
$2$ & $8$ & $6$ & $6=3\cdot2$ \\
\bottomrule
\end{tabular}
\end{center}
In each case $x^3\equiv r^3\equiv r\equiv x\pmod3$.
The table is exhaustive because every integer has one of these residues.
\sourcecredit{Example 3.4.4}
\end{frame}

\begin{frame}[fragile]{Lean: congruence transport and numerical verification}
\begin{lstlisting}
-- In the textbook's modular-arithmetic environment:
example {a b : Int}
    (ha : Int.ModEq 7 a 3) (hb : Int.ModEq 7 b 4) :
    Int.ModEq 7 (a ^ 2 + b) (3 ^ 2 + 4) := by
  rel [ha, hb]

example : Int.ModEq 7 13 6 := by
  use 1
  numbers
\end{lstlisting}
The chapter file also imports \texttt{Library.Tactic.ModEq}.

For residue cases, \texttt{mod\_cases hx : x \% 3} creates the three cases used on the preceding slide.
\sourcecredit{Section 3.4 and its companion Lean file}
\end{frame}

\begin{frame}{Practice 3.4: substitution, witnesses, and cases}
\begin{enumerate}
\item If $n\equiv2\pmod5$, show $n^2+3n+1\equiv1\pmod5$.
\item Find an integer $k$ such that $4k\equiv3\pmod7$.
\item Show that every integer square is congruent to $0$ or $1$ modulo $4$.
\end{enumerate}
\vspace{1em}
Use a different method in each problem: substitute, choose a witness, and split into residues.
\end{frame}

\begin{frame}{Practice 3.4: solution discussion}
\textbf{1.} Substitute the known residue:
\[
n^2+3n+1\equiv2^2+3\cdot2+1=11\equiv1\pmod5.
\]
\textbf{2.} Take $k=6$, since $4\cdot6-3=21=7\cdot3$.
\medskip

\textbf{3.} For residues $r=0,1,2,3$ modulo $4$:
\[
r^2=0,1,4,9\quad\text{has residue}\quad0,1,0,1.
\]
The power rule transfers this calculation from $r$ back to any integer having residue $r$.
\end{frame}

\section{3.5 Bezout's identity}
\begin{frame}{3.5\quad Expressing $1$ as an integer combination}
The integers $3$ and $7$ satisfy
\[
(-2)\cdot3+1\cdot7=1.
\]
Multiplying by any integer $n$ gives
\[
n=-2(3n)+7n.
\]
\begin{block}{Why this helps}
If $3n$ is a multiple of $7$, then both terms on the right are multiples of $7$.
Their sum $n$ must therefore also be a multiple of $7$.
\end{block}
A relation expressing $1$ as an integer combination of two numbers is a B\'ezout identity for those numbers.
\sourcecredit{Section 3.5; method of Example 3.5.1}
\end{frame}

\begin{frame}{Worked example: justify a cancellation by an explicit witness}
Assume $n\in\Z$ and $7\mid3n$. Choose $t\in\Z$ with $3n=7t$.
Then
\[
\begin{aligned}
n&=-2(3n)+7n&&(-2)\cdot3+7=1\\
 &=-2(7t)+7n&&\text{the divisibility witness}\\
 &=7(n-2t)&&\text{algebra}.
\end{aligned}
\]
Hence $7\mid n$, with witness $n-2t$.
\begin{block}{Compare with the failed cancellation modulo $6$}
Here the conclusion follows from a specific integer identity.
Nonzero factors alone do not justify cancellation in congruences.
\end{block}
\sourcecredit{Section 3.5; numerical variant of Example 3.5.2}
\end{frame}

\begin{frame}{The general cancellation statement behind the example}
\begin{block}{Theorem}
Let $m,n,u,v,k\in\Z$, and suppose $um+vn=1$.
If $m\mid nk$, then $m\mid k$.
\end{block}
\textbf{Proof.} Choose $t\in\Z$ with $nk=mt$. Then
\[
\begin{aligned}
k&=(um+vn)k\\
 &=m(uk)+v(nk)\\
 &=m(uk)+v(mt)\\
 &=m(uk+vt).
\end{aligned}
\]
The conclusion has witness $uk+vt$.
\sourcecredit{Generalization of the witness method in Section 3.5}
\end{frame}

\begin{frame}{Worked example: combine two divisibility assumptions}
Suppose $5\mid z$ and $7\mid z$. Write $z=5a$ and $z=7b$.
Since $3\cdot5-2\cdot7=1$,
\[
\begin{aligned}
z&=15z-14z\\
 &=15(7b)-14(5a)\\
 &=35(3b-2a).
\end{aligned}
\]
Thus $35\mid z$.
\begin{block}{Check the arithmetic relationship first}
From $4\mid12$ and $6\mid12$, we cannot conclude $24\mid12$.
The product conclusion needs more than two divisibility assumptions.
\end{block}
\sourcecredit{Section 3.5; numerical variant of Example 3.5.3}
\end{frame}

\begin{frame}{Practice 3.5: find an identity equal to $1$}
\begin{enumerate}
\item Prove that $6\mid5n$ implies $6\mid n$, for every integer $n$.
\item If $4\mid z$ and $9\mid z$, prove that $36\mid z$.
\end{enumerate}
\vspace{1em}
Useful identities to investigate:
\[
6-5=1,\qquad9-2\cdot4=1.
\]
State the new divisibility witness explicitly.
\end{frame}

\begin{frame}{Practice 3.5: solution discussion}
\textbf{1.} Choose $t\in\Z$ with $5n=6t$. Then
\[
n=6n-5n=6n-6t=6(n-t).
\]
The new witness is $n-t$.
\medskip

\textbf{2.} Write $z=4a$ and $z=9b$. Then
\[
z=9z-8z=9(4a)-8(9b)=36(a-2b).
\]
The new witness is $a-2b$.
\end{frame}

\begin{frame}{Chapter 3: definitions turn number theory into algebra}
\begin{center}
\begin{tabular}{ll}
\toprule
Statement & What you must exhibit \\
\midrule
$n$ is even & $n=2k$, with $k\in\Z$ \\
$n$ is odd & $n=2k+1$, with $k\in\Z$ \\
$a\mid b$ & $b=ak$, with $k\in\Z$ \\
$a\equiv b\pmod m$ & $a-b=mk$, with $k\in\Z$ \\
\bottomrule
\end{tabular}
\end{center}
\begin{itemize}
\item Use existing witnesses to build the witness required by the conclusion.
\item Congruence respects addition, subtraction, multiplication, and natural powers.
\item Cancellation needs a justified arithmetic condition; a B\'ezout identity supplies one.
\end{itemize}
\end{frame}

\begin{frame}{Reading and source references}
Heather Macbeth, \emph{The Mechanics of Proof}, Chapter 3:
\begin{itemize}
\item 3.1 Definitions; parity
\item 3.2 Divisibility
\item 3.3 Modular arithmetic: theory
\item 3.4 Modular arithmetic: calculations
\item 3.5 B\'ezout's identity
\end{itemize}
\small
\url{https://hrmacbeth.github.io/math2001/03_Parity_and_Divisibility.html}
\par\medskip
Companion code: \url{https://github.com/hrmacbeth/math2001}
\par\medskip
These slides independently present the definitions, proofs, and new numerical exercises.
The power-rule induction is developed in Chapter 6.
\end{frame}
\end{document}
